\setlength{\unitlength}{0.1mm}
\begin{figure}

\begin{picture}(1370,650)(0,-650)

\put(   0, -650){\framebox(1370,650)[c]{}}

\put( 270, -160){\vector(3,-4){150}}
\put( 270, -260){\vector(3,-2){150}}
\put( 270, -360){\vector(1, 0){150}}
\put( 270, -460){\vector(3, 2){150}}
\put( 270, -560){\vector(3, 4){150}}

\put( 820, -360){\vector(1,0){120}}

\put( 420, -380){\vector(-3, 4){150}}
\put( 420, -380){\vector(-3, 2){150}}
\put( 420, -380){\vector(-1, 0){150}}
\put( 420, -380){\vector(-3,-2){150}}
\put( 420, -380){\vector(-3,-4){150}}

\put( 940, -380){\vector(-1,0){120}}

\put(  10,  -80){\makebox(310,80)[c]{带有原生数据的}}
\put(  10, -110){\makebox(310,80)[c]{处理器}}
\put(  60, -210){\framebox(210,80)[c]{{\F 1}}}
\put(  60, -310){\framebox(210,80)[c]{{\F 2}}}
\put(  60, -410){\framebox(210,80)[c]{{\F 3}}}
\put(  60, -510){\framebox(210,80)[c]{{\F \ldots}}}
\put(  60, -610){\framebox(210,80)[c]{{\F naproc}}}

\put( 370,  -80){\makebox(500,80)[c]{含单个谱分量的}}
\put( 370, -120){\makebox(500,80)[c]{一维全网格数组}}
\put( 420, -400){\framebox(400,60)[c]{活动}}
\put( 420, -460){\dashbox{5}(400,240)[c]{}}
\put( 430, -280){\makebox(380,60)[l]{缓冲空间}}

\put( 970,  -80){\makebox(270,80)[c]{对应的}}
\put( 970, -120){\makebox(270,80)[c]{二维数组}}
\put( 940, -520){\framebox(330,300)[c]{传播}}

\put( 370, -570){\dashbox{15}(950,400)[c]{}}
\put( 400, -580){\makebox(270,80)[l]{\small 在目标处理器上}}

\end{picture}

\caption{分布式内存模式版本中的数据转置。首先，在汇集操作中，数据在图中
从左向右移动。计算完成后，数据在分散操作中从右向左移动。}
\label{fig:transpose}

\botline
\end{figure}
